\section{What the paired Hecke algebra does not control}
\label{sec:paired-algebra-selmer-limit}

\begin{proposition}[AC-G43: a fixed paired algebra with arbitrary cyclic cohomology depth]
\label{AC-G43}
Put $O=\mathbb Z_3[\omega]$, $\lambda=\omega-1$, and
$A=\operatorname{Frac}(O)/O$.  For every $N\geq1$ there
are a finite group $\Gamma_N$, a cubic character $\chi$,
an integral two-dimensional representation $V$, and a
subgroup $D\subset\Gamma_N$ with the following properties.
The reduction of $V$ is absolutely irreducible even on
$\ker\chi$, and there is $\tau$ with $\chi(\tau)=\omega$
and $V(\tau)=1$.  The algebra of paired characteristic
coefficients of $V$ and $\chi V$ is, for every $N$,
\[
                     T'=O\times_{O/\lambda O}O.          \tag{G43.1}
\]
Nevertheless, with $A(\chi)$ denoting the indicated
action on $A$,
\[
 \ker\!\left(H^1(\Gamma_N,A(\chi))
                      \longrightarrow H^1(D,A(\chi))\right)
                  \simeq O/\lambda^N O.                 \tag{G43.2}
\]
Thus its length is arbitrary, even though it is finite,
cyclic and has a one-dimensional socle, while the paired
algebra and its branch separation cost are unchanged.
This is a finite-group cohomology fragment.  It is not
a number-field Selmer example or a countermodel to all
fields and hypotheses of the source \texttt{NData}.
\end{proposition}
\begin{proof}
Let $M_N=O/\lambda^N O$ additively and let $C_3=\langle
\tau\rangle$ act on $M_N$ by multiplication by $\omega$.
Set
\[
 \Gamma_N=(M_N\rtimes C_3)\times Q_8,
 \qquad D=C_3\times Q_8,
 \qquad \chi(m,\tau^a,q)=\omega^a.
\]
Choose $u\in\mathbb Z_3$ with $u^2=-2$ and $u\equiv1\pmod3$.
It exists by lifting the root $1$ modulo $3$: at each
step the correction modulo the next power of $3$ is
uniquely solvable because $2u$ is a unit.  The matrices
\[
 I=\begin{pmatrix}0&1\\-1&0\end{pmatrix},\qquad
 J=\begin{pmatrix}1&u\\u&-1\end{pmatrix}
\]
satisfy $I^2=J^2=-1$ and $IJ=-JI$, and give the required
representation of $Q_8$.  Inflate it to $\Gamma_N$.
Over the algebraic closure of the residue field, $I$
has two distinct eigenlines with eigenvalues $i,-i$,
and the invertible matrix $J$ interchanges them.
There is no common invariant line.  This proves
absolute residual irreducibility, including on
$\ker\chi$, since that subgroup still contains $Q_8$.
The asserted common element is $\tau$ itself.

Let $T'$ include scalar diagonal elements and be
generated by all paired characteristic coefficients.
The coefficient of degree $j$ in the twisted
representation is multiplied by $\chi(g)^j$, which
is congruent to $1$ modulo $\lambda$.  Hence $T'$ is
contained in the right side of (G43.1).  At $g=\tau$
the paired traces are $(2,2\omega)$, and $2$ is a
unit.  Thus $(1,\omega)\in T'$, so $(0,\lambda)$
and, by subtraction from a diagonal element,
$(\lambda,0)$ belong to $T'$.  Every pair with
congruent entries modulo $\lambda$ is a diagonal
element plus a multiple of $(0,\lambda)$.
This proves (G43.1) directly.

The group $Q_8$ acts trivially on $A(\chi)$ and has
order prime to $3$, so every homomorphism from $Q_8$
to $A$ is zero.  A cocycle on $C_3$ is determined by
its value at $\tau$; the cocycle constraint imposes
no additional condition because $1+\omega+\omega^2=0$.
Every such value is a coboundary, because
$\lambda A=A$.  Consequently $H^1(D,A(\chi))=0$.

Restriction of a cocycle on $\Gamma_N$ to $M_N$ is
additive.  After subtracting a coboundary, its
restriction to the split subgroup $D$ is zero.
The cocycle identity then says that its restriction
$f:M_N\to A$ satisfies $f(\omega m)=\omega f(m)$;
it is therefore $O$-linear.  Conversely, every such
map gives the cocycle
\[
                    c_f(m,\tau^a,q)=f(m).
\]
Indeed multiplication in the semidirect product gives
$f(m+\omega^a m')=f(m)+\omega^a f(m')$, exactly the
cocycle identity.  Coboundaries vanish on $M_N$, so
distinct $f$ give distinct classes.  It follows that
\[
 H^1(\Gamma_N,A(\chi))
   =\operatorname{Hom}_O(O/\lambda^N O,A)
   \simeq A[\lambda^N]\simeq O/\lambda^N O.
\]
Since restriction to $D$ has zero target, this is
(G43.2).  Its socle is $\lambda^{N-1}O/\lambda^N O$,
one-dimensional over $O/\lambda O$, as claimed.
\end{proof}

\paragraph{Source direction and scope.}
The algebra in (G43.1) satisfies the common-inertia
and residual hypotheses isolated in AC-G39 and has
the same normalized extra branch loss $1/2$ for
every $N$.  Formula (G43.2) therefore rules out an
upper bound for a general strict cohomology depth
from that algebra alone, even after adding finite
cyclicity and a one-dimensional socle.  The source's
arithmetic local conditions and its actual maps
from congruence modules to cohomology carry additional
information.  This example does not satisfy or
refute those unconstructed maps.  In particular it
does not reverse \texttt{LinkN} or provide a lower
bound for an arithmetic Selmer group.  It is a
directionality check, not closure of a main gate.
